\scanpage{24}\mspage{6 - 42 -}\hfill n° 392

2) Les hypothèses sur $A$ et $B$ étant les mêmes que dans l'exerc. 1, soit $C$ un anneau tel que $A\subset C\subset B$ et que le corps des fractions de $C$ soit égal à $L$ ; tout idéal $\ne0$ dans $C$ est alors un \uline{réseau} de $L$, (relativement à $A$).

a) Montrer que si $\underline a$, $\underline b$ sont deux idéaux de $C$ tels que $\underline a$ soit pseudo-inversible, on a $\chi(\underline a\underline b,C)=\chi(\underline a,C)+\chi(\underline b,C)$.

b) On dit que deux idéaux $\underline a$, $\underline b$ de $C$ sont \uline{pseudo-étrangers} si, pour tout idéal premier $\underline p$ de hauteur 1 dans $A$, $\underline a_{\underline p}$ et $\underline b_{\underline p}$ sont étrangers dans $C_{\underline p}$. Montrer que si $\underline a$ et $\underline b$ sont pseudo-étrangers et $\ne0$, on a $\chi(\underline a\underline b,C)=\chi(\underline a,C)+\chi(\underline b,C)$. (Se ramener au cas où $A$ est un anneau de valuation discrète, observer que l'on a $\unsure{\chi|\chi'}(\underline a,\underline a+\underline b)=\chi(\underline a\cap\underline b,\underline b)$ et utiliser la relation (\quad) du n° \quad du \S{} 4).

c) Déduire de b) que si $\underline a$ est un idéal $\ne0$ de $C$, et si $\underline q_i$, $(1\leq i\leq r)$, sont les idéaux primaires de la décomposition primaire réduite de $\underline a$, correspondant aux idéaux premiers non immergés, on a $\chi(\underline a,C)=\sum_i\chi(\underline q_i,C)$.

d) Soient $\underline q$ un idéal primaire de $C$, $\underline p$ l'idéal premier correspondant. Si $\underline p$ n'est pas de hauteur 1, on a $\chi(\underline q,C)=0$ ; dans le cas contraire, si $\underline p_{\unsure{o|0}}=\underline p\cap A$, on a $\chi(\underline q,C)=\operatorname{deg}_A(\underline q).\underline p_{\unsure{o|0}}$.

e) Soit $\underline a$ un idéal de $C$ ; pour que $\underline a$ soit pseudo-inversible, (exerc. 1 b)), il faut et il suffit que l'on ait $\chi(\underline a,C)=\chi(\underline a B,B)$. (Même méthode que dans l'exerc. 1 b)).

3) Les hypothèses sur $A$, $B$ et $C$ sont les mêmes que dans l'exerc. 2. On dit encore que le réseau $C^*$ est la \uline{codifférente} de $C$, et que $C:C^*$, qui est un idéal de $C$, est la \uline{différente} de $C$ sur $A$ ; on la note $\underline D_{C/A}$.

a) Montrer que le \uline{conducteur} $\underline f_{C/B}$ de $B$ dans $C$, (chap. V, \S{} 1, n° \quad), est égal à $B^*:C^*$, et que pour tout idéal oremier $\underline p$ de hauteur 1 dans $A$, on a $(\underline f_{C/B})_{\underline p}C^*_{\underline p}=B^*_{\underline p}$, (utiliser le fait que $B_{\underline p}$ est un anneau principal).

b) Montrer que l'on a $\underline D_{C/A}\subset\underline f_{C/B}\cap\underline D_{B/A}$, (utiliser a)).

c) Montrer que, pour que $C^*$, considéré comme idéal fractionnaire de $C$, soit pseudo-inversible, (exerc. 1 b)), il faut et il suffit que pour tout idéal premier $\underline p$ de hauteur 1 dans $A$, on ait $B_{\underline p}(\underline D_{C/A})_{\underline p}=(f_{C/B})_{\underline p}(\underline D_{B/A})_{\underline p}$, (utiliser a)).

d) Si $C^*$ est pseudo-inversible, tout idéal fractionnaire $\underline a\ne0$ de $C$ tel que $\underline a:\underline a=C$ est pseudo-inversible, (utiliser le fait que pour deux idéaux fractionnaires non nuls $\underline a$, $\underline b$ de $C$, on a $(\underline a\underline b)^*=\underline a^*:\underline b$).
% archive sheet number: 23
